\section{Experimental design}\label{sec:design}
\subsection{Synthetic laws and matched information comparisons}
All primary evidence is synthetic, so the harmful prevalence and full joint law are known. At the finest resolution $K^*=8$, conditional band probabilities are
\begin{equation}\label{eq:law}
 w_k(\lambda)=\frac{\exp\{\lambda(k-(K^*+1)/2)\}}
 {\sum_{j=1}^{K^*}\exp\{\lambda(j-(K^*+1)/2)\}},\qquad
 \Pbb(J^*=k\mid L=\ell)=(1-q_\ell)w_k(\lambda_\ell),
\end{equation}
with $\Pbb(J^*=\infty\mid L=\ell)=q_\ell$ and $L\sim\operatorname{Bernoulli}(\theta)$. Positive $\lambda$ shifts resolved mass later. Analyses at $K=1,2,4,8$ merge adjacent bands of the same finest law; they do not regenerate a different timing distribution. Terminal age is fixed at eight synthetic time units.

\begin{table}[!htbp]\centering
\caption{Principal generating laws. The strong-timing law and matched control share $\theta,q_1,q_0$, hence identical endpoint masses. Additional validity and boundary laws are specified in the Supplement.}\label{tab:laws}
\begin{tabular}{lrrrrr}\toprule
Law & $\theta$ & $q_1$ & $q_0$ & $\lambda_1$ & $\lambda_0$\\\midrule
Timing and terminal, harmful late & .05 & .30 & .05 & .45 & -.15\\
Strong timing & .02 & .30 & .20 & 1.50 & -.60\\
Matched no timing & .02 & .30 & .20 & 0 & 0\\
Low prevalence & .01 & .30 & .05 & .45 & -.15\\
Independent resolution & .05 & .10 & .10 & 0 & 0\\\bottomrule
\end{tabular}
\end{table}

This decision problem spans 6 laws, 15 sensitivity budgets, 4 partitions and 6 risk budgets. There are thus 90 law--sensitivity pairs, 360 law--sensitivity--partition profiles, and 2,160 decision instances. For each matching legal value of the production upper endpoint, we compare to an oracle obtained by directly analysing the total information in each cell of the contingency table (to 100 decimal places), finding the oracle's optimum avoiding the production endpoint profile, and bounding the range of achievable upper endpoints. These checks numerical correctness of the calculation, and are complementary to the theoretical argument for sharpness above.

\subsection{Power of pre-specified evaluation}
We considered 13 test configurations, leading to 32 pairs of configuration--partition, each replicated by 100 sequences of events. The 'strong' timing law was chosen via a grid search over the population space for large timing gain over its corresponding matched null law (whereby the matching removes only differences in the timing rates). All configuration parameters (laws, sensitivity and risk budgets, partitions and decision timepoints) were logged locally at 19th Oct 2026 prior to confirming the results below; no external pre-registration was done. We reserved stream indices below 1e6 for exploration and 1e6 through 1e6+99 for confirmatory results. Note that for any pair of decisions within a partition and law, we reuse the finest granularity stream of events and so they are dependent: we cannot treat such pairs as independent experimental replicates.

Some exploratory analyses guided our choice of time horizons and choice of comparison procedures ahead of these confirmatory evaluations. Such analyses used separate indices of streams to the above confirmatory 100 per pair and are not combined below. Examples include setting the 'strong' timing horizon to 50,000 events, or restricting our evaluation of certain comparison procedures due to computation requirements ('matched control'). Certifying frequencies below use only streams from the confirmed settings.

As each event matures, we record our confidence sequence over categories and whether it excludes the underlying truth (the data-generating observable law), and take decisions at a set of predefined times. Timepoints of 'strong timing' and 'matched controls' are $\{2000,5000,10000,20000,50000\}$, with 'cannot certify at end' $\{10000,20000,50000\}$ and 'low-prevalence' $\{500, 1000, 2000, 5000, 10000, 20000\}$. Other pairs of decision rules have timepoints summarised and final decision outcomes in the Supplement. All event counts are for matured events, and not wall-clock time. Each partition takes fixed risk $\delta=0.05$ and the decision to accept/reject a class, and margin, of the population space is made ex ante to simulating streams.

Our key performance metric is the empirical frequency that each of 100 streams will (ever) certify, along with a 95\% confidence interval (Clopper--Pearson, \citep{clopper1934}) across the 100 sequences. This quantifies MC noise in estimating the certifying frequency and should not be confused with an interval estimate of $\theta$ at a timepoint. Time of initial certification is the first decision time achieving this (administratively censored if only achieved at the final timepoint), with failed streams remaining. Sequences ceasing updates via empty-set will not contribute subsequent certifications but remain counted in the (ever) certification rate. Our target was for the LCL of the ever-certification frequency to reach 0.8 at the last decision time for 'pop.'-certifiable settings. Note there is no deconvolution of the non-independence induced by multiple decision times in any inferential quantity. All inference is per-condition, with no adjustment for multiple pairs.

\subsection{Validation and Comparisons} 
A further study across 7 settings was done at a shorter horizon with more streams (500 sequences of length 500 events, decision times every 100 events) to validate the implementation of the procedure and its risk-projection, including cases where the number of bands exceeds 32 and we fill in by minimum information. Empirical 'ever violated' frequency (one-sided binomial UCL @ 95\%) is compared against 0.06 as an empirical analog of the configured risk tolerance. Target false-alarm is 0.05. Results at this shorter timepoint are only informative on validity, and not certification at longer decision times.

The main competitor we report against is an end-point decision by the same construction (\emph{i.e.,} $\rho$, budget of risk, stream/time-uniform) using only the final confidence sequence over trajectory endpoint. We also compare to 2 other procedures that address related-but-different problems. One repeats static category-Wilson (Bonferroni) inference after each maturing event and optimizes for risk similarly to our problem (using the same sensitivity function) but is not guaranteed to be valid as a time-uniform procedure and is only approx- valid under the static distribution at each timepoint. Another is a delay-ignorable procedure (Bernoulli-mixture confidence sequence) applied to the resolved labels and requires stronger assumptions; we expect this can apply on the $q_1=q_0, \lambda_1 = \lambda_0$ class of laws only. Behaviour outside these classes is diagnostic but cannot be construed as out-performing the other (valid) competitors.

We validate our numeric computations with oracles (above), theorem equivalence and an independent verification of feasibility of our stored risk bounds. We scale up to K up to 128 (i.e. $1,2,4,\ldots,128$), warmup x 5 times then time over x30 invocations at each K-value on detached workers (restricted to one-thread operation of relevant numerics packages). For each K, we regenerate the required instance of the conditional distribution with +0.01 nats of known information and do not time those. As such, this does not evaluate successive refinement. Timing on the outer target at each K includes the generation and application of a 500 length balanced prefix (deterministic) and projection. All code is written in Python [Numpy, SciPy, mpmath, python-flint] and reproducible with locked dependencies available on the code repo. We did not capture enough detail to state performance portably.

\subsection{Evidence provenance}\label{sec:provenance} 
All decision-times of certifying, and associated confidence intervals, are aggregated from 3,200 (power) and 3,500 (validation) single-stream evaluations. Dependency manifests do not encode a commit hash of our local working directory, therefore we are unable to trace which implementation gave rise to which past evaluations. Details given in the Supplement. 